\scanpage{64}
% block: numerical characters item 3
\handwritten{3}) La donnée de l'application
\[
\handwritten{\sigma\unsure{\mathrel{\vert\! -\! -}|\mapsto}}(Z_{\handwritten{\sigma}},\handwritten{\xi_\sigma})\ ,\ \handwritten{\xi_\sigma\in}\operatorname{Br}(Z_{\handwritten{\sigma}})
\]
envisagée dans a).

% block: numerical characters item 4
4) La donnée de la fonction
\[
\operatorname{rang}\ :\ \handwritten{\mathcal{M}\longrightarrow}\underline{Z}
\]
lorsque $\handwritten{\mathcal{M}}$ est supposée\deleted{semi-}tannakienne.\deleted{\unsure{(i.e.}}

% block: effective part
On notera que la donnée de 1),2) équivaut à la donnée de l'anneau commutatif $K=K(\handwritten{\mathcal{M}})$, avec sa ``partie effective
\[
\begin{aligned}
K^{\mathrm{eff}}&=\text{sous-groupe de }K(\handwritten{\mathcal{M}})\text{ engendré par les}\\
&\qquad c\handwritten{\ell}(M),\ M\handwritten{\in}\operatorname{Ob}\quad,\\
&=\underline{N}^{(\handwritten{\Sigma})}
\end{aligned}
\]
(NB $\handwritten{\Sigma}$ se récupère à partir de \deleted{\unsure{ces}}\inserted{\handwritten{la}} donnée\inserted{\handwritten{$K^{\mathrm{eff}}\subset K$}} comme l'ensemble des éléments non nuls minimaux de $K^{\mathrm{eff}}$). La donnée de 3) permet alors de trouver, comme explicité dans a), la variance fonctorielle de l'anneau $K(\handwritten{\mathcal{M}}_{k'})$ par rapport à $k'$. Enfin la donnée de 4) permet en principe de retrouver les polynômes caractéristiques (en particulier, la trace et le déterminant) d'un endomorphisme d'un objet \deleted{de}(semi-simple ?) $M$ de $E$, en termes de polynômes caractéristiques réduits (resp. trace réduite, norme réduite) dans les algèbres d'Azumaya\inserted{$\operatorname{End}(M_i)$,} correspondant aux composantes isotypiques $M_i$ de $M$. Pour $M$ isotypique de rang $n$, on aura en effet, si $K=\operatorname{End}(M)$ a comme centre $Z$ de rang $r$ sur $k$, donc $\unsure{m|n}=n'\unsure{\deleted{d}r|\deleted{r}d|r}$, et $K$ de rang $d^2$ sur $Z$, avec $d\handwritten{\mid}n'$ (NB $M$ de rang $n'$ sur $Z$):
\[
\begin{aligned}
\det_M f&=\deleted{\det}N_{Z/k}(\operatorname{Nr}_{K/Z}(f))^{n'/d}\\
\operatorname{Tr}_M f&=\operatorname{Tr}_{Z/k}(\operatorname{Trr}_{K/Z}(f))\ \frac{n'}{d}\\
\operatorname{Pol}_M(f,t)&=N_{Z/k}(\operatorname{Polr}_{K/Z}(f,t))^{n'/d}\quad\inserted{(n'/d=n/dr)}\ .
\end{aligned}
\]
\margin{\handwritten{\unsure{Signaler} que la connaissance des car. num. de $\mathcal{M}$ implique celle des car. numériques de $\mathcal{M}_{k'}$ pour \unsure{toute} extension $k'$ de $k$, (``fonctorielle en $k'$'')}}

% block: warning about reconstruction
On fera attention que la connaissance des caractères numériques de la catégorie tannakienne $\handwritten{\mathcal{M}}$ ne permet pas en général de réconstituer celle-ci à équivalence près, même si $k$ est un corps alg. clos de car. 0 (de sorte que la donnée de 3) devient vide) et $\handwritten{\mathcal{M}}$ semi-simple
% archive sheet number: 63
